\noindent\textbf{Elección del procedimiento.} La velocidad terminal corresponde a aceleración nula, no velocidad nula. El equilibrio vertical es $mg=\rho V_t^2AC_D/2$; dividir por el coeficiente positivo y tomar raíz permite calcular cada altitud.

\begin{enumerate}

\item Usar primero la densidad de nivel del mar para la comparación.

\[\fitmath{\mathrm{Vsea}=\sqrt{\frac{2\cdot m\cdot g}{rho0\cdot A\cdot C_D}}}\]

\[\fitmath{\mathrm{Vsea}=\sqrt{\frac{2\cdot (80)\cdot (9.81)}{(1.225)\cdot (0.7)\cdot (1)}}\approx 42.78361\,\mathrm{m/s}}\]

\item Repetir con la densidad a 3000 m.

\[\fitmath{\mathrm{Vt}=\sqrt{\frac{2\cdot m\cdot g}{\rho\cdot A\cdot C_D}}}\]

\[\fitmath{\mathrm{Vt}=\sqrt{\frac{2\cdot (80)\cdot (9.81)}{(0.9091)\cdot (0.7)\cdot (1)}}\approx 49.66377\,\mathrm{m/s}}\]

\end{enumerate}
